\section{原子操作的延迟与吞吐}
\label{l13:sec:throughput}

\subsection{同一地址上的串行服务}

吞吐率（throughput）表示单位时间完成多少次操作，延迟（latency）表示一次操作从发起到完成经历多久。对普通内存访问，多个独立请求可以重叠，因此一个请求延迟较长，并不必然使总吞吐率等于该延迟的倒数。原子热点的特殊之处在于，很多请求依赖同一个地址上前一次更新的结果。

考虑一个简化模型：同一地址每次原子更新独占服务时间为 $\tau$，下一次冲突更新必须在前一次服务完成后才能开始。对持续饱和的请求序列，这个地址的最高更新速率约为
\begin{equation}
R_{\mathrm{addr}}\le\frac{1}{\tau}.
\label{l13:eq:atomic-rate}
\end{equation}
这里 $\tau$ 是模型中的持续服务间隔，不是某个线程排队后测到的总等待时间。若同时到达 $q$ 个更新，最后一个请求可能需要等待约 $q\tau$ 才完成，但设备在这段时间内仍以每 $\tau$ 完成一次更新的速率工作。

假设这个服务过程需要超过 $1000$ 个时钟周期，那么同一热点地址每周期完成的更新少于 $1/1000$ 次。相对于“每周期一次更新”的理想基准，速率降到了不足千分之一。这个比较的对象是单个热点地址的更新能力；若多个地址可以独立服务，它们仍可能并行贡献总吞吐。

令桶 $j$ 一共接收 $H_j$ 次更新。若采用相同服务时间 $\tau$，即使忽略读取输入、其他桶和所有调度开销，完成统计的时间仍受到最长串行链约束：$T_{\mathrm{hot}}\ge\max_j H_j\tau$。所有输入进入同一个桶时，这条链长为 $N$；输入接近均匀分布时，每个桶大约只承担 $N/B$ 次更新。这个下界说明了输入分布如何影响竞争，但多个桶仍可能共享内部执行资源，因此它本身不是完整的设备性能预测。

更多线程有助于提出更多独立请求，但不会让同一地址上必须依次生效的更新同时完成。如果原子热点已经饱和，继续增加等待线程主要会增加排队长度。优化可以从两方面着手：缩短服务时间，或者减少需要到达同一个公共地址的更新数量。

\subsection{从显存到二级缓存和共享内存}

若原子读—改—写必须经过动态随机存取存储器（DRAM），一次更新需要取得旧值并提交新值，长距离数据传输和存储访问会使完整服务时间较大。在以完整读、写往返计时的模型中，这包含读延迟、写延迟以及内部路由和传输开销。对同地址冲突操作，后续更新还必须等待当前更新的完整结果。

二级缓存（L2 cache）可以缩短这一路径。目标位置驻留在 L2 时，支持相应原子处理的硬件可以在更近的位置完成更新，避免完整的 DRAM 往返。不同线程块仍然访问同一公共全局对象；如果它们都更新同一桶，冲突仍存在，只是每次服务可能更快。

共享内存（shared memory）属于线程块，存储距离更近，原子更新通常具有更短的服务路径。它不能让所有线程块直接访问同一个共享数组。要用共享内存统计整份输入，需要让各块各自维护局部直方图，随后把局部结果合并到全局输出。因此，从 L2 的硬件加速进一步走向共享内存，需要算法配合。

\begin{table}[htbp]
\centering
\caption{三种原子更新位置及其作用}
\label{l13:tab:memory-levels}
\small
\begin{tabularx}{\textwidth}{>{\raggedright\arraybackslash}p{2.0cm}YY}
\toprule
更新位置 & 延迟来源与改进 & 仍需处理的限制 \\
\midrule
DRAM 路径 & 读、写往返与内部传输构成较长服务过程 & 同地址排队，完整服务时间较长 \\
L2 路径 & 目标驻留缓存时可避免 DRAM 往返 & 全部线程块仍可能争用同一公共桶 \\
共享内存 & 块内就近更新局部桶 & 块内同桶更新仍串行，需要跨块合并 \\
\bottomrule
\end{tabularx}
\end{table}

这里的层次比较说明机制，不为某一设备指定固定延迟。实际性能还依赖原子数据类型、目标分布、缓存状态和存储系统实现。无论在哪个层次，同一个计数器仍必须表现为一串正确的原子更新。

\subsection{性能优化的两个方向}

缩短服务时间时，更新次数可以完全不变。例如，把一个块的 $n_r$ 次更新转移到共享内存后，这些输入仍各贡献一次加一，但每次操作经过较短的局部路径。减少公共更新次数时，则先把多次相同目标的贡献合成局部计数，再向公共桶提交一个较大的增量。

设某线程块处理的输入中有 $400$ 个值等于 $7$。直接算法会对公共桶 $7$ 做 $400$ 次原子加一；局部统计可以先在该块中得到计数 $400$，最后只对公共桶 $7$ 做一次原子加 $400$。每个输入的贡献仍然保留，而公共地址只接收一次更新。共享内存私有化同时利用了更近的局部存储和较少的公共更新。

评价这种变换时，要计入新增的局部数组清零、阶段同步和结果合并。它是否划算取决于每个块能把多少输入更新聚合成多少个非零桶。下一节以完整算法和操作数量明确这些条件。
